\section{Data and methods}
\label{sec:methods}

This section defines the three surface fields computed throughout the paper,
fixes the sign conventions that make them reproducible, describes the
polyhedral solver and the shape models, specifies the mesh operators used to
probe resolution behaviour, and reports the verification of the implementation
against exact analytic solutions.

\subsection{Target quantities}
\label{sec:targets}

Let the body occupy a region $B$ bounded by a closed triangulated surface with
$\Nf$ facets, and let the interior density $\dens$ be uniform. The
gravitational potential at a field point $\mathbf{r}$ is

\begin{equation}
U_{g}(\mathbf{r}) \;=\; -\,G\dens \int_{B}
\frac{\mathrm{d}^{3}r'}{\lvert \mathbf{r}-\mathbf{r}' \rvert},
\label{eq:Ug}
\end{equation}

so that $U_{g}<0$ and the self-gravitational attraction is

\begin{equation}
\ngrav(\mathbf{r}) \;=\; -\nabla U_{g}(\mathbf{r}),
\label{eq:g}
\end{equation}

which points into the body. For a body rotating uniformly at angular rate
$\omegarot$ about the $z$ axis of the body-fixed frame, write
$\mathbf{r}_{\perp}=(x,y,0)$ for the component perpendicular to the spin axis.
The centrifugal acceleration is directed away from the spin axis,

\begin{equation}
\acf(\mathbf{r}) \;=\; \omegarot^{2}\,\mathbf{r}_{\perp},
\label{eq:acf}
\end{equation}

and the effective acceleration experienced by a particle at rest in the
rotating frame is

\begin{equation}
\geff(\mathbf{r}) \;=\; \ngrav(\mathbf{r}) + \acf(\mathbf{r}).
\label{eq:geff}
\end{equation}

The \emph{effective gravitational slope} $\slope$---hereafter simply the
gravitational slope, the qualifier being understood---of facet $i$ with outward
unit normal $\nhat_{i}$ is the angle between the upward direction $-\geff$ and
that normal,

\begin{equation}
\slope_{i} \;=\; \arccos\!\left(
\frac{-\,\geff(\mathbf{c}_{i})\cdot\nhat_{i}}
     {\lVert \geff(\mathbf{c}_{i}) \rVert} \right),
\label{eq:slope}
\end{equation}

evaluated at the facet centroid $\mathbf{c}_{i}$. A surface perpendicular to
the local effective gravity has $\slope=0$; a surface on which the effective
gravity has an outward component has $\slope>90^{\circ}$, which is possible
only where the centrifugal term dominates or where the local topography
overhangs.

The \emph{geopotential} adds the rotational term to the gravitational
potential,

\begin{equation}
\Upot(\mathbf{r}) \;=\; U_{g}(\mathbf{r})
\;-\;\tfrac{1}{2}\,\omegarot^{2}\,\lVert\mathbf{r}_{\perp}\rVert^{2},
\label{eq:geopot}
\end{equation}

so that $\geff=-\nabla\Upot$ and the minima of $\Upot$ on the surface mark the
sinks toward which loose material is driven.

Two scalar summaries of these fields are used throughout. Weighting each facet
by its area $A_{i}$, the area-weighted mean slope is

\begin{equation}
\slopeavg \;=\;
\frac{\sum_{i} A_{i}\,\slope_{i}}{\sum_{i} A_{i}},
\label{eq:slopeavg}
\end{equation}

and the area-weighted geopotential mean and variance are

\begin{equation}
\Ubar \;=\; \frac{\sum_{i} A_{i}\,\Upot_{i}}{\sum_{i} A_{i}},
\qquad
\potvar^{2} \;=\;
\frac{\sum_{i} A_{i}\,\bigl(\Upot_{i}-\Ubar\bigr)^{2}}{\sum_{i} A_{i}}.
\label{eq:potvar}
\end{equation}

The square root $\potvar$ of Eq.~\eqref{eq:potvar} is the corresponding
area-weighted standard deviation. Because $\Upot$ carries units and its zero
point is fixed by Eq.~\eqref{eq:Ug}, comparisons across bodies use the
dimensionless \emph{normalized geopotential dispersion}

\begin{equation}
\potvarn \;=\; \frac{\potvar}{\lvert \Ubar \rvert}.
\label{eq:potvarn}
\end{equation}

which is the coefficient of variation of the surface geopotential. We use the
term \emph{dispersion} rather than \emph{variance} throughout, because
$\potvarn$ is built from the standard deviation and not from
$\potvar^{2}$.

Both $\slopeavg$ and $\potvarn$ are area-weighted rather than facet-weighted,
so that neither is biased by the uneven facet areas that decimation produces.

\subsubsection{Sign conventions}
\label{sec:signs}

Equations~\eqref{eq:g}--\eqref{eq:geopot} fix every sign, but two conventions
in common use differ from them and are worth stating explicitly, because a
discrepancy in either is easy to introduce and hard to detect.

First, some solvers---including the library used here---return the positive
Newtonian integral $-U_{g}$ rather than $U_{g}$ itself, together with an
acceleration vector that points into the body. The implementation negates the
returned potential before forming Eq.~\eqref{eq:geopot}, so that the
gravitational and rotational terms of the geopotential carry consistent signs.

Second, the centrifugal term must \emph{oppose} self-gravity, as in
Eq.~\eqref{eq:acf}: at the equator of a rotating body the effective
acceleration is smaller in magnitude than the attraction alone. The opposite
sign is easy to adopt by mistake when the underlying solver stores gravity as
an outward-pointing vector, and---as Section~\ref{sec:numval} shows---it cannot
be detected by testing a non-rotating body.

\subsubsection{Surface-limit convention}
\label{sec:surflimit}

The Newtonian potential and its gradient possess well-defined surface limits
for a homogeneous solid of finite density carrying no surface mass layer: the
physical acceleration is finite and, away from edges and vertices, continuous
across a face. In a polyhedral implementation, however, evaluating exactly on a
face makes the solid-angle and edge terms of the analytic expression ambiguous
at the level of a branch choice. Fields are therefore evaluated at

\begin{equation}
\mathbf{p}_{i} \;=\; \mathbf{c}_{i} \;-\; \epsilon\,\nhat_{i},
\qquad
\epsilon \;=\; 10^{-4}\,L,
\label{eq:eps}
\end{equation}

where $L$ is the diagonal of the shape model's bounding box, so that the
evaluation point lies just inside the surface and the interior limit of the
numerical expression is selected consistently. The displacement is four orders
of magnitude below the body scale and two to three orders below a typical facet
dimension, and the results are insensitive to it: halving or doubling
$\epsilon$ changes $\slopeavg$ and $\potvarn$ in the fifth significant figure.

\subsection{Polyhedral field evaluation}
\label{sec:solver}

The self-gravitation of a homogeneous polyhedron is computed with the analytic
line-and-face-integral method of \citet{werner1996}, in the formulation of
\citet{tsoulis2012}, as implemented in the open-source \texttt{polyhedral-%
gravity} package \citep{gravitylib}. The method is exact for a uniform
polyhedron---no series truncation, no mascon approximation---and is valid for
non-convex and multiply concave surfaces, which is essential for the bilobed and
irregular bodies considered here. Facet normals are oriented outward before
evaluation, and the solver is run in its outward-normal mode.

Since Eqs.~\eqref{eq:Ug}--\eqref{eq:g} are linear in $\dens$, the field for any
assumed density follows from a single unit-density evaluation by scaling. This
is used to sweep the assumed bulk density at fixed shape
(Section~\ref{sec:validation}) at the cost of one solver call, and it motivates the
region-superposition strategy discussed as a future interior-density extension
in Section~\ref{sec:interior}.

All lengths in the shape files are in kilometres; internally they are converted
to SI, and slopes are reported in degrees.

\subsection{Shape models}
\label{sec:shapes}

Table~\ref{tab:shapes} lists the nine shape models used, covering seven
asteroids and spanning oblate top shapes, elongated and bilobed configurations,
angular and irregular fragments, and one extreme dog-bone. Two bodies are
represented by two independent models each, which permits the cross-version
comparison of Section~\ref{sec:crossversion}.

Bulk densities are the measured values where a spacecraft or a satellite has
constrained the mass, and representative taxonomic values otherwise; because
the fields scale with $\dens$ in a known way, an assumed density affects the
absolute slope but not the resolution behaviour that is the subject of
Sections~\ref{sec:slopesens} and~\ref{sec:varrobust}. For Ryugu the density is
set to $1200\ \mathrm{kg\,m^{-3}}$ exactly, matching the value used for the
reference products against which we validate, rather than to the marginally
different measured bulk value.

\begin{table}[t]
\centering
\small
\caption{Shape models used in this work. $\Nf$ is the facet count of the
delivered model. Densities are in \si{\gram\per\cubic\centi\metre} and
rotation periods in hours.}
\label{tab:shapes}
\begin{tabular}{llrrrl}
\toprule
Body & Model & $\Nf$ & $\dens$ & $P$ & Role \\
\midrule
(101955) Bennu     & OLA v20            & 196\,608 & 1.194 & 4.2961  & validation \\
(162173) Ryugu     & SFM 49k v20180804  &  49\,152 & 1.200 & 7.627   & validation \\
(433) Eros         & Gaskell            &  49\,152 & 2.670 & 5.270   & resolution \\
(25143) Itokawa    & Gaskell            &  49\,152 & 1.950 & 12.132  & resolution, cross-version \\
(25143) Itokawa    & radar-based        &   3\,688 & 1.950 & 12.132  & cross-version \\
(951) Gaspra       & Thomas             &  32\,040 & 2.700 & 7.042   & resolution \\
(243) Ida          & Thomas             &  32\,040 & 2.600 & 4.634   & resolution \\
(216) Kleopatra    & Ostro radar        &   4\,092 & 3.380 & 5.385   & cross-version \\
(216) Kleopatra    & \texttt{kleo.v2.7} &   2\,292 & 3.380 & 5.385   & cross-version \\
\bottomrule
\end{tabular}
\end{table}

Every model is checked before use. The mesh is verified to be watertight and
manifold, with no degenerate facets and no inward-pointing normals; the
principal axes of inertia are computed and the spin axis confirmed to coincide
with the axis of maximum inertia, so that the latitude used in the profiles of
Section~\ref{sec:validation} is meaningful; and the bounding box is compared
with the published dimensions of the body, which is the check that catches unit
errors and mis-scaled models. This last test is not a formality: it is what
identifies the two Kleopatra models as differing in scale as well as in detail
(Section~\ref{sec:crossversion}).

\subsection{Mesh operators}
\label{sec:operators}

Two operators are applied to the meshes, and they are deliberately different in
character.

\paragraph{Decimation.} To vary resolution while preserving the body, vertices
are clustered on a uniform grid: the bounding box is divided into cells, all
vertices falling in a cell are replaced by their centroid, faces that become
degenerate are removed, and duplicate faces are collapsed. The grid is refined
until the target vertex count is reached. This reduces $\Nf$ while leaving the
overall figure of the body essentially unchanged.

It is important that decimation does \emph{not} smooth the surface. A coarser
facet spans the same topography that several finer facets spanned before, so
its normal still reflects the local relief; what changes is how that relief is
sampled, not whether it is present. This distinction is what separates the
resolution effect of Section~\ref{sec:slopesens} from the roughness effect
below.

\paragraph{Smoothing.} To remove facet-scale relief while holding $\Nf$ fixed,
Laplacian smoothing is applied: each vertex is displaced toward the centroid of
its topological neighbours,

\begin{equation}
\mathbf{v}_{j} \;\leftarrow\; \mathbf{v}_{j}
+ \lambda\left( \frac{1}{\lvert N(j)\rvert}\sum_{m\in N(j)}\mathbf{v}_{m}
- \mathbf{v}_{j}\right),
\qquad \lambda = 0.5,
\label{eq:laplace}
\end{equation}

iterated a prescribed number of times. Successive iterations progressively
suppress short-wavelength relief while leaving the global figure of the body
intact. Applying the two operators separately makes it possible to attribute a
change in $\slopeavg$ either to how the surface is sampled (decimation) or to
how rough it is (smoothing).

\subsection{Resolution diagnostics}
\label{sec:diagnostics}

For a scalar diagnostic $X$ evaluated on a set of meshes with facet counts
$\Nf^{(1)}<\dots<\Nf^{(K)}$ obtained by decimating a single model, two measures
are reported.

The \emph{range sensitivity} is the peak-to-peak variation normalized by the
mean over the set,

\begin{equation}
S[X] \;=\; \frac{\max_{k} X^{(k)} - \min_{k} X^{(k)}}
{\tfrac{1}{K}\sum_{k} X^{(k)}}.
\label{eq:sens}
\end{equation}

The \emph{convergence residual} is the fractional change between the two finest
meshes,

\begin{equation}
C[X] \;=\;
\frac{\bigl\lvert X^{(K)} - X^{(K-1)} \bigr\rvert}{X^{(K-1)}}.
\label{eq:conv}
\end{equation}

Both are quoted as percentages. By construction $C[X]\ge 0$: it records the
\emph{magnitude} of the fractional change over the final refinement step. Where
the direction of that change is also of interest, the corresponding signed
quantity

\begin{equation}
C_{\mathrm{signed}}[X] \;=\;
\frac{X^{(K)} - X^{(K-1)}}{X^{(K-1)}}
\label{eq:convsigned}
\end{equation}

may be reported alongside it. Every sequence studied here increases
monotonically over its final refinement step, so the two coincide in magnitude
and $C_{\mathrm{signed}}$ is positive throughout.

The two answer different questions. $S[X]$ measures how much a reported value
would change had a coarser or finer mesh been used anywhere in the tested range,
and is the quantity relevant when comparing published numbers. $C[X]$ measures
whether the diagnostic is approaching a limit as the mesh is refined, and is
the quantity relevant to asking whether a resolution-independent value exists
at all. A diagnostic may have a modest $S$ merely because it varies slowly, yet
a large $C$ because it has not settled; the two are reported together
throughout.

\subsection{Numerical verification}
\label{sec:numval}

The implementation is verified against two closed-form solutions, chosen so
that between them they exercise every term in
Eqs.~\eqref{eq:Ug}--\eqref{eq:geopot}.

\paragraph{Potential: homogeneous ellipsoid.} For a uniform ellipsoid with
semi-axes $a\ge b\ge c$ the exterior potential is given exactly by the
classical elliptic integral

\begin{equation}
U_{g}(\mathbf{r}) \;=\; -\pi G\dens\,abc \int_{\lambda}^{\infty}
\left(1 - \frac{x^{2}}{a^{2}+u} - \frac{y^{2}}{b^{2}+u}
        - \frac{z^{2}}{c^{2}+u}\right)
\frac{\mathrm{d}u}{\Delta(u)},
\label{eq:ellipsoid}
\end{equation}

with $\Delta(u)=\sqrt{(a^{2}+u)(b^{2}+u)(c^{2}+u)}$ and $\lambda$ the largest
root of the ellipsoidal coordinate equation, which vanishes on the surface. In
the limit $a=b=c=R$ this reduces to $-GM/R$ on the surface, and the quadrature
used reproduces that limit exactly. Evaluated on a triangulated ellipsoid with
$a,b,c=1500,1000,800$~m and $\dens=2000\ \mathrm{kg\,m^{-3}}$, the polyhedral
potential agrees with Eq.~\eqref{eq:ellipsoid} facet by facet to a Pearson
correlation of $1.000000$, with a mean relative error of $0.18\%$ and a maximum
of $0.21\%$; the residual is consistent with the polygonal approximation of the
curved surface.

\paragraph{Slope: uniformly rotating sphere.} For a homogeneous sphere of
radius $R$ and surface gravity $g=GM/R^{2}$ rotating at rate $\omegarot$, write
$k=\omegarot^{2}R/g$ for the ratio of centrifugal to gravitational acceleration
at the equator. Combining Eqs.~\eqref{eq:acf}--\eqref{eq:slope} gives the slope
at latitude $\phi$ in closed form,

\begin{equation}
\cos\slope(\phi) \;=\;
\frac{1-k\cos^{2}\phi}
     {\sqrt{(1-k)^{2}\cos^{2}\phi+\sin^{2}\phi}}.
\label{eq:rotsphere}
\end{equation}

At $k=0.5$ this predicts $19.11^{\circ}$, $18.43^{\circ}$ and $13.90^{\circ}$
at latitudes $30^{\circ}$, $45^{\circ}$ and $60^{\circ}$; the implementation
returns $18.97^{\circ}$, $18.33^{\circ}$ and $13.91^{\circ}$, agreeing to
better than $0.15^{\circ}$ at the mesh resolution used.

This second test is not redundant, and the reason is worth stating. Reversing
the sign of $\acf$ in Eq.~\eqref{eq:geff} is equivalent to replacing $k$ by
$-k$ in Eq.~\eqref{eq:rotsphere}, which at $k=0.5$ changes the predicted slope
at $45^{\circ}$ latitude from $18.43^{\circ}$ to $11.31^{\circ}$---a
discrepancy of nearly two-fifths that the rotating test detects immediately.
A \emph{non-rotating} sphere, by contrast, returns zero slope for either sign,
and so cannot detect the error at all. Since a static sphere is the natural
first test to write, we recommend that codes of this kind be checked against
Eq.~\eqref{eq:rotsphere} as well; both tests are provided with the code.

As a control, the non-rotating sphere is also run and returns a mean slope of
$0.35^{\circ}$, consistent with the polygonal discretisation of a curved
surface at the resolution used.

\subsection{Reproducibility}
\label{sec:repro}

All computations use Python~3.13 with \texttt{numpy} and \texttt{scipy}, and
version~3.3.1 of the \texttt{polyhedral-gravity} package \citep{gravitylib}
for the self-gravitation. No random numbers enter the analysis: every quantity
reported is deterministic given the shape model, the assumed density, and the
rotation period.

The analysis code, the verification scripts implementing
Eqs.~\eqref{eq:ellipsoid} and~\eqref{eq:rotsphere}, the mesh diagnostics, and a
machine-readable file containing every numerical result quoted in this paper are
archived and openly available; sources for all shape models are listed there.
Shape models themselves are not redistributed.
